%======================================================================
\section{An Example R Markdown Document}
%======================================================================

\sectionframe

\begin{frame}{About This Example}
\begin{itemize}
  \item This section reproduces the knitted HTML output of the default RMarkdown template (\code{R\_Intro\_Markdown.Rmd}) created in the previous section, with a few equations added.
  \item Compare it with your own knitted file: the text, the R code chunk, its output, and the plot all come from the same \code{.Rmd} file.
\end{itemize}
\end{frame}

\begin{frame}[fragile]{R Markdown}
\begin{itemize}
  \item This is an R Markdown document. Markdown is a simple formatting syntax for authoring HTML, PDF, and MS Word documents. For more details on using R Markdown see \url{http://rmarkdown.rstudio.com}.
  \item When you click the \textbf{Knit} button a document will be generated that includes both content as well as the output of any embedded R code chunks within the document. You can embed an R code chunk like this:
\end{itemize}
\begin{lstlisting}[language=R]
summary(cars)
\end{lstlisting}
\begin{lstlisting}[style=output]
     speed           dist
 Min.   : 4.0   Min.   :  2.00
 1st Qu.:12.0   1st Qu.: 26.00
 Median :15.0   Median : 36.00
 Mean   :15.4   Mean   : 42.98
 3rd Qu.:19.0   3rd Qu.: 56.00
 Max.   :25.0   Max.   :120.00
\end{lstlisting}
\end{frame}

\begin{frame}{Including Plots}
\begin{itemize}
  \item You can also embed plots, for example:
\end{itemize}
\begin{center}
\includegraphics[width=0.45\textwidth]{ch04/rmd_pressure_plot}
\end{center}
\begin{itemize}
  \item Note that the \code{echo = FALSE} parameter was added to the code chunk to prevent printing of the R code that generated the plot.
  \item You can also write equations (LaTeX syntax), e.g.
  \[
    y_i = X_i \beta + \epsilon_i, \qquad X = U D V^T
  \]
\end{itemize}
\end{frame}
